%% Computational Geosciences — Springer Nature LaTeX submission
%% Template: sn-jnl v3.1, option [iicol,sn-mathphys-num]
%% All figure files at root level for submission (Fig1.png–Fig4.png).

\documentclass[pdflatex,sn-mathphys,Numbered,iicol]{sn-jnl}

%% ---- Packages ----
\usepackage{graphicx}
\usepackage{multirow}
\usepackage{amsmath,amssymb,amsfonts}
\usepackage[title]{appendix}
\usepackage{xcolor}
\usepackage{textcomp}
\usepackage{booktabs}
\usepackage{manyfoot}
\usepackage{float}
\usepackage{hyperref}
\usepackage{cleveref}

%% ---- Custom commands ----
\newcommand{\logK}{\ln K}
\newcommand{\bxi}{\boldsymbol{\xi}}
\newcommand{\bxiast}{\boldsymbol{\xi}^{*}}
\newcommand{\bxitrue}{\boldsymbol{\xi}_{\mathrm{true}}}
\newcommand{\nobs}{n_{\mathrm{obs}}}
\newcommand{\kappaJ}{\kappa(\mathbf{J})}
\newcommand{\kappaHGN}{\kappa(\mathbf{H}_{\mathrm{GN}})}

%% ---- Manuscript info ----
\title[Separating observability and workflow limitations]{Separating linearized observability, direct recoverability, and workflow limitations in physics-informed neural network inversion}

\author*[1]{\fnm{Fei} \sur{Yang}}\email{yangf@cigem.cn}

\affil[1]{\orgname{China Institute of Geo-Environment Monitoring}, \orgdiv{China Geological Survey}, \street{Minzu University South Road}, \city{Beijing}, \postcode{100081}, \country{China}}

\abstract{Physics-informed neural networks (PINNs) are increasingly used for groundwater inverse problems, but a good hydraulic-head fit does not establish conductivity recovery. We present a controlled diagnostic study of PINN-KLE conductivity inversion for steady two-dimensional groundwater flow, aiming to separate local information content, gradient transfer, round-selection quality, and workflow-limited recovery---where recoverable information exists in the direct reduced map but is not reliably transmitted by the staged workflow.

Gradient audits show that the head-observation loss provides no direct gradient pathway to KLE coefficients; updates are mediated through the PDE-residual branch in the implemented staged workflow, whose direction is weakly aligned with truth. A no-truth round selector identifies the best alternating round in 17/18 cases, yet only 3/18 weak-pass conductivity recovery. The direct FDM-KLE baseline weak-passes 24/24 cases across eight truth realizations ($\sigma_{\log K}=0.17$--$0.80$), confirming that the sparse inverse problem contains recoverable low-dimensional information; the failure is therefore workflow-limited relative to the matched direct baseline. Jacobian diagnostics show full linearized local rank for 3/5/10-mode direct maps at $\nobs=200$ and for 3-mode maps across 50--400 observation points, with conditioning burden increasing with mode count. Local sensitivity at hybrid-recovered coefficients remains comparable to true-point conditioning (condition-number ratio 0.919--1.010, median 0.953, across twelve paired cases), indicating that failures do not arise from loss of local direct-map observability. The paper argues for a more careful diagnostic standard: head fitting, linearized local observability, direct deterministic baselines, and closed-loop conductivity recovery must be reported separately.}

\keywords{Groundwater inverse problem, Physics-informed neural networks, Karhunen--Lo\`eve expansion, Local observability, Workflow diagnostics}

\pacs[MSC Classification]{65N21, 86A05, 68T07, 35R30, 65N06}

%% =====================================================================
\begin{document}
\maketitle

%% =====================================================================
\section{Introduction}
\label{sec:intro}

Groundwater conductivity inversion is a classical ill-posed inverse problem.
Sparse hydraulic-head observations provide indirect and smoothed information
about subsurface heterogeneity, and the answer obtained from an inverse model
depends strongly on parameterization, prior structure, boundary forcing, and
the richness of hydraulic excitation~\cite{mclaughlin1996reassessment,
carrera1986aquifer1, carrera1986aquifer2, kitanidis1995quasilinear,
zhou2014inverse, yeh2000hydraulic}. The identifiability of aquifer
parameters from limited head data has been studied extensively:
Sun and Yeh~\cite{sun1990identifiability} introduced coupled inverse-analysis
and experimental design concepts, Hill and
Tiedeman~\cite{hill2007effective} developed sensitivity-based calibration
methods, and Tarantola~\cite{tarantola2005inverse} and Aster et
al.~\cite{aster2018parameter} provided general inverse-theory frameworks for
parameter estimation. Hydraulic tomography further demonstrated that multiple
pumping tests can improve subsurface
characterization~\cite{zhu2005characterization, cardiff2013hydraulic,
yeh2015flow}. Physics-informed neural networks offer an attractive
differentiable framework for combining neural approximators with partial
differential equations~\cite{raissi2019pinn, karniadakis2021piml}, and they
have been applied to subsurface-flow parameter learning and groundwater
inverse modeling~\cite{tartakovsky2020subsurface, guo2023htpinn,
bandai2022forward, ji2024groundwater, gong2025klepinnc}. However, PINNs do
not remove the information-content problem that makes groundwater inversion
difficult. A neural model can match observed heads while still failing to
identify conductivity parameters.

This paper focuses on that evaluation gap. In some recent PINN inverse
studies, the reported evidence is dominated by head misfit, PDE residual, or
visual agreement. For conductivity inversion studies, these diagnostics are
insufficient unless they are connected to coefficient-level recovery. In a
PINN-KLE setting, one must ask whether the data term produces gradients with
respect to the KLE coefficients, whether alternating training rounds can be
selected without using full-field truth, whether recovered coefficients pass an
internal consistency check, and whether a direct physics-based inverse
baseline can recover coefficients in cases where the tested PINN workflow does
not.

We study these questions in a deliberately controlled steady groundwater
benchmark. The benchmark is not intended to reproduce all features of a field
aquifer. Instead, it is designed to make different failure modes separable. A
failure may be structural, meaning that the observations and equations do not
determine the parameter. It may be a practical recovery failure under finite
noisy data. It may be local uncertainty in a shallow objective valley. Or it
may be workflow-limited: a direct baseline can recover the reduced parameters,
but the tested PINN-hybrid workflow cannot transmit that information into a
conductivity estimate.

The contributions are:
\begin{enumerate}
\item a controlled numerical diagnostic study for sparse-head PINN-KLE
  groundwater inversion;
\item a loss-structure gradient audit showing that the standard head-only data
  loss does not directly update KLE coefficients;
\item a no-truth round-selection audit that separates choosing the best
  available alternating round from achieving conductivity recovery;
\item a matched direct FDM-KLE least-squares baseline showing that the main
  moderate-noise 3-mode failure is consistent with a workflow-limited
  interpretation in the tested setting;
\item sparse-head Jacobian, Fisher-information, and Gauss-Newton diagnostics
  showing linearized local observability for the direct map, increasing
  conditioning burden with mode count, and comparable conditioning at
  hybrid-recovered coefficients---supporting the interpretation that the workflow failure does not
  arise from loss of local direct-map observability at the recovered parameters.
\end{enumerate}

The central claim is deliberately narrow: the tested PINN-hybrid workflow can
lose coefficient information that is still accessible to a direct deterministic
FDM-KLE inverse baseline. The paper does not claim that all PINN inverse
formulations fail, nor that the numerical thresholds found here transfer
unchanged to real aquifers.

%% =====================================================================
\section{Problem setup and parameterization}
\label{sec:setup}

\subsection{Governing equation and boundary conditions}
\label{sec:pde}

All experiments solve a steady two-dimensional saturated groundwater flow
problem on the unit square:
\begin{equation}
  \nabla \cdot \bigl( K \, \nabla h \bigr) = 0, \quad (x,y) \in (0,1)^2, \quad K(x,y) > 0.
  \label{eq:pde}
\end{equation}
The default boundary conditions impose fixed heads at the left and right
boundaries and no-flow conditions at the top and bottom:
\begin{equation}
  h\big|_{x=0} = 1, \quad h\big|_{x=1} = 0, \quad
  \mathbf{n} \cdot K \nabla h \,\big|_{y=0,\,y=1} = 0.
  \label{eq:bc}
\end{equation}

This source-free steady setup is intentionally simple. It removes pumping,
recharge, transient response, irregular geometry, and conceptual-model
uncertainty so that workflow-level failures can be interpreted without field
ambiguity. The limitation is equally important: numerical thresholds reported
here are benchmark-specific and should not be treated as monitoring-network
design rules for real aquifers.

\subsection{KLE log-conductivity model}
\label{sec:kle}

The heterogeneous log-conductivity field is represented with a truncated KLE:
\begin{equation}
  \logK(x, y) = \overline{\logK}
    + \sum_{i=1}^{M} \xi_i \, \sqrt{\lambda_i} \, \phi_i(x, y),
  \label{eq:kle}
\end{equation}
where $\overline{\logK}$ is the mean log-conductivity, $\lambda_i$ and
$\phi_i$ are the eigenvalues and eigenfunctions of the spatial covariance
kernel, $\xi_i$ are the dimensionless stochastic coefficients, and $M$ is
the truncation order. The coefficient vector
$\bxi = (\xi_1, \ldots, \xi_M)^\top$ is the target reduced-order coefficient
vector of the inverse problem.
Most matched-noise comparisons use a 3-mode KLE target ($M=3$). Additional
mode-bridge experiments examine $M=5$ and $M=10$ settings. The KLE
representation follows standard KLE formulations for random fields and flow in
heterogeneous media~\cite{ghanem1991stochastic, zhang2004efficient,
lu2004conditional}. It is used as a controlled reduced parameterization, not as
a universal geological model. It lets us ask whether low-dimensional
information is present in the data and whether the tested workflow can recover
it.

\subsection{Synthetic observations and closed-loop validation}
\label{sec:obs}

Synthetic sparse head observations are sampled from FDM reference solutions
and perturbed by controlled Gaussian noise. Recovery quality is evaluated
through a closed-loop FDM solve: the estimated conductivity coefficients are
inserted back into the FDM forward model, and the resulting hydraulic head is
compared with the reference field. The main metrics are:
\begin{itemize}
\item closed-loop head RMSE;
\item log-conductivity RMSE;
\item KLE coefficient correlation;
\item weak-pass and strong-pass operational recovery labels (weak-pass:
      $\log K$ RMSE $<0.15$ and $h$ RMSE $<0.006$; strong-pass: $\log K$
      RMSE $<0.10$ and $h$ RMSE $<0.003$).
\end{itemize}
The pass labels are audit thresholds, not physical phase transitions. They are
used to summarize whether a run achieved a useful coefficient recovery under
the stated benchmark criteria. The closed-loop check uses the same FDM forward
family as the reference solution; it is an internal consistency check on the
coefficient estimate, not an independent model validation against a different
physics solver.

%% =====================================================================
\section{Diagnostic methods}
\label{sec:methods}

\subsection{PINN-KLE gradient audit}
\label{sec:grad-audit}

In the tested staged PINN-KLE formulation, a neural network approximates the head
field, while the KLE coefficients are treated as trainable variables. The loss
contains head-observation, boundary-condition, and PDE-residual terms. We
compute component-wise gradients of these loss terms with respect to the KLE
coefficients. This diagnostic asks a simple question: which terms actually
provide coefficient-level gradient information?

The gradient audit is not a convergence test. It is a structural diagnostic of
the loss pathway. In the tested workflow, the head-observation loss updates only
the neural head approximator and provides no direct gradient pathway to $\bxi$;
therefore, head fitting alone cannot by itself serve as sufficient evidence of conductivity
recovery. This diagnostic is workflow-specific and does not describe all PINN
formulations---it reflects the staged-freeze design of the tested workflow.

\subsection{Hybrid PINN-FDM workflow}
\label{sec:workflow}

The staged PINN-FDM workflow separates head approximation from coefficient
recovery:
\begin{enumerate}
\item Stage A fits a neural head field using sparse observations and physics
  constraints.
\item Stage B freezes the sampled head field on a grid and optimizes KLE
  coefficients through a discrete residual.
\item Stage C performs an FDM forward solve using the recovered coefficients
  and evaluates closed-loop metrics.
\end{enumerate}
Alternating Stage A/B/C rounds are repeated. Earlier experiments used oracle
full-field head RMSE to choose the best round, which is not available in field
applications. The current diagnostic uses sparse-observation RMSE as a
no-truth selector and audits it against oracle head RMSE only in synthetic
mode.

\subsection{Direct FDM-KLE least-squares baseline}
\label{sec:baseline}

The direct baseline optimizes the same KLE coefficients against the same
sparse head observations using repeated FDM forward solves. It is a
deterministic parameter-estimation baseline in the same reduced KLE space.
This baseline is not intended to mimic every feature of PEST, EnKF, or other
hydrogeologic inversion systems. Its role is sharper: if the direct FDM-KLE
baseline recovers the coefficients while the PINN-hybrid workflow fails, then
the failure cannot be attributed simply to absence of low-dimensional
information in the observations.

\subsection{Workflow-diagnostics aggregation}
\label{sec:diag-aggregate}

The workflow diagnostics package combines five post-processing audit layers.
Each layer addresses a specific failure-mode question. None of the layers
requires re-running the forward or inverse model; they aggregate existing
training outputs into a diagnostic chain.

\paragraph{Layer 1: Gradient-transfer audit.}
For each loss component $\mathcal{L}_k$ (head observation, boundary condition,
PDE residual), we record the gradient vector
$\nabla_{\bxi} \mathcal{L}_k$ at selected training steps. The audit asks:
which loss terms produce non-zero gradients with respect to the KLE
coefficients? If the data loss provides no direct computational gradient
pathway, then head fitting alone cannot drive coefficient updates. We also compute the cosine similarity
between the PDE-loss gradient direction and the direction from the current
$\bxi$ to $\bxitrue$ to assess gradient alignment.

\paragraph{Layer 2: Selector-versus-recovery audit.}
After all alternating rounds are completed, we evaluate two independent
criteria: (a) the \emph{selector} quality, defined as the sparse-observation
RMSE of the chosen round relative to the oracle-best round, and (b) the
\emph{recovery} quality, defined as the closed-loop logK RMSE of the chosen
round. A round can have excellent selector quality (near-zero oracle regret)
but poor recovery quality. This separation is essential because field
applications must select rounds without access to truth.

\paragraph{Layer 3: Matched direct-baseline gap.}
For each PINN-hybrid case, we run a direct FDM-KLE least-squares inversion
(Section~\ref{sec:baseline}) using the same truth, same observation
locations, and same noise realization. The gap between the direct baseline's
logK RMSE and the hybrid's logK RMSE measures workflow-limited information
loss. If the direct baseline recovers the coefficients and the hybrid does
not, the failure is not attributable to absence of low-dimensional
information in the data.

\paragraph{Layer 4: Late-round degradation audit.}
We track the per-round Stage A head RMSE, Stage B logK RMSE, and Stage C
closed-loop head RMSE across alternating rounds. A transition is classified as
\emph{degrading} if the Stage C RMSE worsens from one round to the next. For
each degrading transition, we record whether Stage A head RMSE, Stage B logK
RMSE, or both rebounded, and compute the correlation between Stage B
logK-RMSE change and Stage C head-RMSE change. This identifies whether late
degradation is driven by Stage A fitting failure, Stage B coefficient damage,
or A/B coupling.

\paragraph{Layer 5: Mode-bridge diagnostic.}
We compare hybrid recovery with a perfect-head coefficient recovery reference.
The perfect-head reference solves the inverse problem using the full-field FDM
head solution as input rather than a neural approximation. The ratio of hybrid
logK RMSE to perfect-head logK RMSE quantifies the workflow damage
attributable to the neural head approximation versus the increasing
conditioning burden of higher mode counts.

\subsection{Sparse-head Jacobian, Fisher information, and Gauss-Newton Hessian}
\label{sec:jacobian}

For the direct FDM-KLE observation map, we compute a finite-difference
sensitivity matrix (the Jacobian) at a reference coefficient vector
$\bxiast$:
\begin{equation}
  \begin{aligned}
  J_{ij} &= \frac{\partial h(\mathbf{x}_i;\, \bxiast)}{\partial \xi_j},\\
  i &= 1,\ldots,\nobs, \quad j = 1,\ldots,M,
  \end{aligned}
  \label{eq:jacobian}
\end{equation}
where $\mathbf{x}_i$ are the sparse observation locations. The
finite-difference approximation uses an absolute perturbation
$\delta = 10^{-3}$:
\begin{equation}
  J_{ij} \approx
    \frac{h(\mathbf{x}_i;\, \bxiast + \delta\,\mathbf{e}_j)
          - h(\mathbf{x}_i;\, \bxiast - \delta\,\mathbf{e}_j)}
         {2\,\delta}.
  \label{eq:fd-jacobian}
\end{equation}
From this matrix we form the local Fisher-information approximation under
least-squares Gauss--Newton linearization and the Gauss-Newton Hessian
approximation~\cite{sun1990identifiability, tarantola2005inverse,
aster2018parameter}:
\begin{equation}
  \mathbf{F} = \frac{1}{\sigma^2} \, \mathbf{J}^\top \mathbf{J},
  \qquad
  \mathbf{H}_{\mathrm{GN}} = 2\,\mathbf{J}^\top \mathbf{J},
  \label{eq:fim-hgn}
\end{equation}
where $\sigma^2$ is the observation-noise variance. Under the standard
assumption of iid Gaussian observation errors with variance $\sigma^2$,
$\mathbf{F}$ corresponds to the Fisher information matrix for the linearized
observation model. Because $\mathbf{F}$ and
$\mathbf{H}_{\mathrm{GN}}$ share the same rank and their singular-value
ratios are determined by $\mathbf{J}$ alone (up to the noise-dependent
scaling $1/\sigma^2$), the Jacobian rank and singular-value spectrum suffice
to diagnose local observability and conditioning of the direct reduced inverse
map. Under the present least-squares Gauss-Newton formulation, scanning the
noise level $\sigma$ only rescales the eigenvalues and does
not change the rank or condition number of $\mathbf{F}$. We note that the
condition numbers satisfy $\kappa(\mathbf{H}_{\mathrm{GN}}) =
\kappa(\mathbf{F})$ and, for the standard spectral condition number,
$\kappa(\mathbf{H}_{\mathrm{GN}}) \approx \kappa(\mathbf{J})^2$; the larger
HGN values reported in the results reflect this squaring relationship and do
not indicate a qualitatively different conditioning regime.

These diagnostics are local and linearization-based: they do not prove global
identifiability, uniqueness, or absence of multimodal
structure~\cite{sun1990identifiability, hill2007effective}. They do, however,
test whether small perturbations in all coefficient directions are detectable by
the sparse-head map at the linearization point. We apply the Jacobian diagnostic at two types of
reference points: (a) the true coefficient vector and (b) the
hybrid-recovered coefficient vector, to test whether workflow
failure coincides with loss of local observability.

%% =====================================================================
\section{Results}
\label{sec:results}

\subsection{Head-observation loss does not directly identify KLE coefficients}
\label{sec:res-grad}

The gradient audit shows that in the tested staged workflow, the
head-observation loss provides no direct gradient pathway to KLE coefficients
both before and after Stage A. In
the representative diagnostic case, the initial PDE-loss coefficient-gradient
norm is $4.84 \times 10^{3}$, while the after-Stage-A PDE-loss
coefficient-gradient norm falls to $1.68$. The PDE gradient direction is only
weakly aligned with the direction to the true coefficients: the cosine is
$0.200$ initially and $0.309$ after Stage A. The first actual
coefficient-update direction has cosine $0.056$ with the direction to truth.

This explains why head fitting alone is a weak inverse-success diagnostic in
the tested PINN-KLE formulation. Coefficient information reaches $\bxi$
through the PDE-residual branch in the implemented staged workflow, not directly
through the sparse head data term. This is consistent with broader PINN evidence that optimization
pathologies and gradient imbalance can strongly affect whether
physics-informed losses transmit useful parameter
information~\cite{wang2021gradient, krishnapriyan2021failure}
(Figure~\ref{fig:workflow}).

\begin{figure*}[t]
  \centering
  \includegraphics[width=\textwidth]{Fig1.png}
  \caption{Workflow diagnostics for the tested PINN-KLE inverse procedure.
    Panel~A summarizes the coefficient-gradient transfer audit. Panel~B
    separates no-truth selector quality from conductivity recovery (fractions
    computed across noise groups of unequal case counts; see Section~4.2).
    Panel~C compares median logK RMSE for direct FDM-KLE and PINN-hybrid cases.
    Panel~D summarizes late-round degradation drivers; bars are not mutually
    exclusive (total transitions~$=$~38, of which a subset co-occur).}
  \label{fig:workflow}
\end{figure*}

\subsection{No-truth round selection is not the same as recovery}
\label{sec:res-selector}

The observation-RMSE selector performs well in the current medium-budget pilot
aggregate, which spans five noise levels ($\sigma=0.005,0.01,0.015,0.02,0.05$)
and therefore extends beyond the compact matrix's three-level scan. Across 18
cases, all 18 are hard-pass selector cases. However,
only 1 of these 18 cases passes the conductivity-recovery criterion. The
median hybrid logK RMSE increases from $0.153$ at noise $\sigma=0.005$ to
$0.178$, $0.236$, $0.318$, and $1.133$ at noise levels $\sigma=0.01$,
$0.015$, $0.02$, and $0.05$, respectively.

The conclusion is not that the selector fails. Rather, it correctly answers a
limited operational question: which available alternating round is best under a
no-truth sparse-observation criterion? That question is distinct from whether
the best available round contains a useful conductivity estimate.

\subsection{Direct FDM-KLE exposes workflow-limited recovery loss}
\label{sec:res-baseline}

The direct FDM-KLE comparison is the central workflow diagnostic. We present
three progressively controlled experimental phases: an initial medium-budget
pilot aggregate, a compact multi-truth matrix, and a targeted three-seed
audit.

\paragraph{Pilot aggregate.}
Across the initial matched aggregate, direct FDM-KLE weak-passes 15/18 cases,
while PINN-hybrid weak-passes 1/18. The median logK RMSE of the direct
baseline remains below or near the weak-pass threshold through moderate noise,
whereas the PINN-hybrid median logK RMSE is already above threshold at noise
$\sigma=0.005$ and grows with noise.

At noise $\sigma=0.01$ and $0.015$, the direct baseline weak-passes all
matched fine transition cases, while the medium-budget PINN-hybrid workflow
weak-passes none. At noise $\sigma=0.02$, the direct baseline again
weak-passes all matched cases in the medium-budget comparison, while the
PINN-hybrid workflow fails. These results support the workflow-limited
interpretation: the sparse direct inverse problem still contains recoverable
low-dimensional information, but the tested PINN-hybrid workflow does not
reliably transmit it into recovered coefficients.

\paragraph{Compact multi-truth matrix.}
As a controlled multi-truth check, we ran direct FDM-KLE on eight
3-mode truth realizations spanning $\log K$ standard deviations from $0.17$ to
$0.80$, at $\nobs=200$ and three noise levels ($\sigma=0.005,0.01,0.02$) using
a single observation seed. The direct baseline passed all 24 cases ($8 \times 3$)
under both weak and strong criteria. Median closed-loop logK RMSE increased with
noise, from $0.018$ at $\sigma=0.005$ to $0.036$ at $\sigma=0.01$ and $0.072$
at $\sigma=0.02$, remaining well below the weak-pass threshold ($\log K$ RMSE $<0.15$, $h$ RMSE $<0.006$). In the
most challenging display truth ($8$, $\sigma_{\log K}=0.80$), logK RMSE reached
$0.089$ at $\sigma=0.02$ but still satisfied the weak-pass criterion. The full
8-truth results are provided in Online Resource~1. By
contrast, the five additional truths are evaluated as FDM-only baselines;
matched PINN-hybrid comparisons remain with the three original display truths
($1$, $2$, $3$) for which seed-audit hybrid data already exist.

The matched PINN-hybrid expansion was then run for the more favorable
$\nobs=200$ subset of the same three truths and three noise levels. The
contrast is sharp. In the matched $\nobs=200$ cases, direct FDM-KLE passes
all 9 cases under both weak and strong criteria. PINN-hybrid produces hard-pass
no-truth selector outcomes in all 9 cases, with zero oracle head-RMSE regret,
but weak-passes only 2/9 and strong-passes 0/9. By noise, PINN-hybrid
weak-passes 1/3 at $\sigma=0.005$, 1/3 at $\sigma=0.01$, and 0/3 at
$\sigma=0.02$; the corresponding median logK RMSE values are $0.590$,
$0.626$, and $0.777$. By truth, the workflow is near the weak-pass boundary
for truth $1$, but fails all tested noise levels for truths $2$ and $3$.
The compact matched matrix therefore reinforces the separation between direct
recoverability, round-selection quality, and workflow recovery.

\paragraph{Three-seed audit.}
To check whether this contrast is an artifact of a single observation/training
seed, we repeated the compact comparison at $\nobs=200$ for the two bracketing
noise levels ($\sigma=0.005$ and $0.02$) using three seeds ($202$, $303$, and
$404$). Across these 18 seed-audit cases, direct FDM-KLE weak-passes all 18
and strong-passes 15. In the matched PINN-hybrid audit, the selector remains
operationally reliable in 17/18 cases with median oracle head-RMSE regret
equal to zero, but conductivity recovery weak-passes only 3/18 and never
strong-passes. The truth dependence persists: truth $1$ weak-passes 3/6
hybrid cases, whereas truths $2$ and $3$ weak-pass 0/6 each. This
targeted seed audit does not replace a full statistical ensemble, but it
removes the most immediate single-seed explanation for the observed
workflow-limited gap. Seeds 202, 303, and 404 were selected as a targeted
bracket rather than a random draw from a full factorial design; the audit
removes the most immediate single-seed explanation but does not establish
seed-averaged statistics.

\subsection{Late-round degradation is Stage-B-led with frequent coupling}
\label{sec:res-degrade}

The alternating-round degradation audit analyzes 11 runs and 38
round-to-round transitions. Nine transitions worsen the Stage C closed-loop
head RMSE. All 9 of these worsening transitions include a rebound in Stage B
logK RMSE, while 6 also include a rebound in Stage A head RMSE. On worsening
transitions, the correlation between the Stage B logK-RMSE change and the
Stage C head-RMSE change is $0.784$.

This result argues against a simplistic explanation in which the workflow
failure is only a Stage A head-fitting problem. Late degradation is better
described as Stage-B-led with frequent Stage A/B coupling near the stopping
point.

\subsection{Direct-map sensitivity shows linearized local observability with
increasing conditioning burden}
\label{sec:res-obs}

The finite-difference sparse-head Jacobian diagnostic gives full relative rank
for the tested 3/5/10-mode direct maps at $\nobs=200$. The Jacobian condition
number increases with mode count
(Table~\ref{tab:cond}, Figure~\ref{fig:sensitivity}).

\begin{table}[t]
  \caption{Jacobian and Gauss-Newton condition numbers for 3/5/10-mode
    direct FDM-KLE maps at $\nobs=200$.}
  \label{tab:cond}
  \centering
  \begin{tabular}{lccc}
    \toprule
    Mode count & Rank & $\kappaJ$ & $\kappaHGN$ \\
    \midrule
    3  & 3/3 & 10.6 & 112 \\
    5  & 5/5 & 35.6 & $1.27\times10^{3}$ \\
    10 & 10/10 & 54.3 & $2.95\times10^{3}$ \\
    \bottomrule
  \end{tabular}
\end{table}

This result does not prove global identifiability. Full column rank of the
Jacobian at a reference point implies linearized local observability---that is,
small perturbations in all coefficient directions are detectable by the
sparse-head map at the linearization point. Global identifiability, uniqueness,
and optimization convergence are not implied. The tested observation layout
shows linearized local observability through 10 modes. Therefore, higher-mode
hybrid failure should not be described as proven absence of local information.
It is more consistent with workflow damage and increasing conditioning burden.

The same 3-mode diagnostic was then repeated at the three compact-matrix display
truth realizations ($1$, $2$, and $3$) with $\nobs=200$. All three cases
remain full rank ($3/3$). Their Jacobian condition numbers are $10.8$,
$7.53$, and $11.4$, respectively, with Gauss-Newton condition numbers $117$,
$56.7$, and $131$. This truth sweep is important because truth $2$, one of
the severe PINN-hybrid failure cases, is not the worst conditioned direct map;
it is the best conditioned of the three in this local diagnostic. Thus, the
compact truth results strengthen the workflow-limited interpretation: the
tested hybrid failures do not coincide with loss of local direct-map
observability at the true coefficients.

To check whether the full-rank finding is an artifact of the observation
density, we repeated the truth-point FIM diagnostic for the same three display truths
across four observation layouts ($\nobs = 50, 100, 200, 400$) at
$\sigma=0.005$. All 12 cases remain full rank ($3/3$). Truth $2$ retains
the best conditioning across all densities, with Jacobian condition numbers
$6.4$, $6.9$, $7.5$, and $7.8$ at $\nobs=50, 100, 200, 400$,
respectively (compared with $9.1$--$11.0$ for truth $1$ and $10.3$--$11.8$
for truth $3$). The FIM rank is unaffected by reducing observations from 400 to
50. This observation-density sweep (Figure~\ref{fig:fim-density-sweep}) supports the workflow-limited
interpretation: rank deficiency is not observed at any tested density, yet the
staged workflow fails to extract that information into recovered coefficients at
$\nobs=200$.

\begin{figure}[!ht]
\centering
\includegraphics[width=\linewidth]{Fig2.png}
\caption{Jacobian condition number and weakest singular value across
observation densities for the three compact-matrix truth realizations. All
Jacobians remain full rank ($3/3$) from 50 to 400 observation points. Truth
$2$ retains the best conditioning at all densities.}
\label{fig:fim-density-sweep}
\end{figure}

We further evaluated local sensitivity at the selected hybrid-recovered
coefficients for both hard-failure truths ($2$ and $3$), using the six
targeted seed-audit cases each ($\sigma = 0.005/0.02$, seeds $202/303/404$,
$\nobs=200$). Across the twelve paired cases,
recovered-$\bxi$ distances range from $1.54$ to $4.59$, showing that the
hybrid workflow places the coefficients far from the truth. Nevertheless,
recovered-point local direct-map conditioning remains comparable to true-point
conditioning. The recovered/true Jacobian condition-number ratio
$\kappa(\mathbf{J}_{\mathrm{rec}}) / \kappa(\mathbf{J}_{\mathrm{true}})$
ranges from $0.919$ to $1.010$ with median $0.953$, and all twelve
recovered-point Jacobians remain full rank ($3/3$). The recovered/true
weakest-singular-value ratio ranges from $0.930$ to $1.022$ with median
$0.979$.

For truth $2$, recovered-$\bxi$ distances are between $4.13$ and $4.59$,
with condition-number ratios $0.935$--$1.010$. For truth $3$,
recovered-$\bxi$ distances are between $1.54$ and $2.57$, with
condition-number ratios $0.919$--$0.959$. In both hard-failure truths, the
hybrid workflow does not appear to fail by moving into a locally unobservable
region of the direct FDM-KLE map. It moves to wrong coefficient values that
still have locally observable sparse-head sensitivity. This replicated finding
across two truths strengthens the workflow-limited interpretation: the failure
mechanism is not loss of local direct-map observability at the recovered
parameters (Figure~\ref{fig:recovered-fim}).

The close agreement of condition numbers between true and recovered points
(ratio $0.919$--$1.010$) indicates that the local sensitivity spectra at the true
and recovered points are comparable. This observation supports the side-by-side
comparison of true-point and recovered-point local diagnostics, but similar
condition numbers do not, by themselves, imply that the two points lie in the
same attraction basin of the optimization landscape.

\begin{figure}[t]
  \centering
  \includegraphics[width=\linewidth]{Fig3.png}
  \caption{Finite-difference sparse-head Jacobian spectrum diagnostics.
    Panel~A shows the singular value spectrum of the observation Jacobian
    $\mathbf{J}$ for 3-, 5-, and 10-mode KLE cases.
    Panel~B shows the Jacobian condition number $\kappa(\mathbf{J})$ (left
    axis, bars) and the weakest singular value (right axis, line) versus
    mode count ($\nobs=200$). The Fisher information and Gauss-Newton Hessian
    spectra are determined from $\mathbf{J}^T\mathbf{J}$ and
    $2\mathbf{J}^T\mathbf{J}$ respectively and are not plotted separately.}
  \label{fig:sensitivity}
\end{figure}

\begin{figure*}[t]
  \centering
  \includegraphics[width=\textwidth]{Fig4.png}
  \caption{Recovered-point versus true-point local conditioning for two
    hard-failure truths ($2$ and $3$). Panels A/C show Jacobian condition
    numbers; panels B/D show weakest singular values. Each panel pairs
    true and recovered values across six seed-audit cases per truth. The
    dashed line marks equal conditioning. All recovered points remain full
    rank ($3/3$) despite coefficient distances of $1.54$--$4.59$ from truth.}
  \label{fig:recovered-fim}
\end{figure*}

\subsection{Mode bridge separates workflow damage from structural claims}
\label{sec:res-bridge}

Mode-bridge diagnostics compare hybrid recovery with fixed-head coefficient
recovery. At 3 modes, workflow damage is visible but still tolerable in the
representative bridge case: hybrid logK RMSE is $0.125$, while perfect-head
recovery gives logK RMSE $1.72 \times 10^{-4}$. At 5 and 10 modes, the
hybrid workflow fails with logK RMSE $1.219$ and $1.426$, respectively, while
perfect-head recovery remains near exact ($7.42 \times 10^{-4}$ and
$1.35 \times 10^{-3}$).

The conclusion is that 5/10-mode failure in the current hybrid workflow is not
evidence of structural impossibility. It is a workflow-limited failure under
the tested staged approximation and recovery procedure.

%% =====================================================================
\section{Discussion}
\label{sec:discussion}

\subsection{Workflow-limited failure is distinct from non-identifiability}
\label{sec:disc-workflow}

The main result is the separation between local information content and
workflow recovery. The direct FDM-KLE baseline and the sparse-head Jacobian
diagnostic show that low-dimensional coefficient information can be present in
the direct inverse map. The PINN-hybrid workflow can still fail to recover
that information. This is not a semantic distinction: it determines whether one
should conclude that the observations are intrinsically insufficient or that
the tested workflow fails to reliably preserve recoverable coefficient information.

For the present benchmark, the moderate-noise 3-mode failure is best described
as workflow-limited. This phrase should be read narrowly. It means
workflow-limited relative to the tested PINN-hybrid design and the direct
FDM-KLE baseline in the same reduced KLE space. It does not imply that all
PINN-based groundwater inverse formulations fail.

For clarity, we distinguish four failure categories relevant to this benchmark:
(i) \emph{information-limited} failure, where the observations do not contain
coefficient-level information in any reduced map; (ii) \emph{optimization-limited}
failure, where information is present in the direct map but the optimizer cannot
reach it; (iii) \emph{architecture-limited} failure, where the neural
representation cannot approximate the required function space; and
(iv) \emph{stage-coupling} failure, where errors in one training stage propagate
irreversibly into subsequent stages (these categories are not mutually
exclusive). The present workflow-limited finding combines
elements of (ii) and (iv): the direct map contains recoverable information, but
the staged A/B optimization pathway does not reliably transmit it into recovered
coefficients. We note that establishing the precise causal chain from the
stage-coupling mechanism (iv) to the observed recovery gap would require tracing
the optimization trajectory from initial to recovered parameters; the present
diagnostic framework identifies the gap and its stage-of-origin, but does not
fully resolve the mechanistic link between them.

\subsection{Relationship to classical inverse modeling}
\label{sec:disc-classical}

The direct FDM-KLE baseline plays a role similar to classical deterministic
parameter estimation in a reduced parameter
space~\cite{doherty2015pest, zhou2014inverse, hill2007effective}. It is not a
full comparison against PEST, EnKF, EnRML, MCMC, or geostatistical inversion
software~\cite{evensen2003enkf, oliver2008inverse, white2018model}.
However, it is strong enough for the specific diagnostic claim made here: if
direct FDM-KLE least squares recovers the same coefficients from the same
observations while PINN-hybrid fails, then the failure cannot be assigned
solely to low-dimensional information absence.

A broader comparison with PEST-style regularized inversion, adjoint
least-squares, or ensemble methods would be valuable, but it is not required
to establish the workflow-limited finding. Those methods would address a wider
question: which inverse family is most accurate or efficient? The present
paper addresses a narrower diagnostic question: whether the tested PINN-hybrid
workflow loses recoverable information. The comparison is diagnostic rather than
competitive: we test whether recoverable coefficient information is accessible
to a direct solver in the same reduced space, not whether FDM-KLE is a superior
inversion method for practical applications.

\subsection{What the FIM and Hessian diagnostics do and do not prove}
\label{sec:disc-fim}

The Jacobian, Fisher-information, and Gauss-Newton diagnostics are local and
linearization-based~\cite{tarantola2005inverse, aster2018parameter}. A full-rank Jacobian
at a reference point does not prove global uniqueness, rule out multimodality,
or quantify posterior uncertainty under nonlinear effects. It only shows that
small perturbations in the tested coefficient directions are locally visible to
the sparse-head observation map. Conversely, increasing condition number
indicates increasing sensitivity burden, not categorical failure.

This local interpretation is sufficient for the paper's logic. The FIM
diagnostic is used to reject an overly strong structural claim: the 5/10-mode
hybrid failures should not be reported as if the direct observation map had no
local coefficient information. They should instead be presented as workflow
failures under a more difficult conditioning regime.

The recovered-point FIM analysis (Section~\ref{sec:res-obs}) adds an
important nuance to this local interpretation. When the same Jacobian
diagnostic is evaluated at the hybrid-recovered coefficients rather than at
the true coefficients, the direct-map remains full rank with comparable
conditioning (recovered/true condition-number ratio: $0.919$--$1.010$,
median $0.953$, across twelve paired cases from two hard-failure truths).
This means that the hybrid workflow does not fail by drifting into a locally
unobservable region of parameter space. It reaches coefficient values that
remain locally observable under the direct linearized map.

This finding has two implications. First, it further strengthens the
workflow-limited interpretation: the failure mechanism is not loss of local
direct-map observability at the recovered parameters. The local sensitivity
spectra at true and recovered points are comparable; what is lost is the optimization pathway
that connects the data to the correct coefficient values. Second, it
illustrates an important limitation of local observability diagnostics: a
full-rank Jacobian at a point does not guarantee that an iterative method can
reach that neighborhood from a distant starting point, nor does it guarantee
that a different optimization landscape (here, the PINN-hybrid loss surface)
shares the same favorable geometry as the direct least-squares landscape.

The distinction between local observability at a point and recoverability by a
specific optimization procedure is therefore not merely semantic. A full-rank
FIM at the true coefficients tells us that the observations contain
low-dimensional information; a full-rank FIM at the recovered coefficients
tells us that the optimization has not collapsed into an unobservable region;
but neither guarantees that the PINN-hybrid training trajectory can connect the
starting point to the correct solution. This three-way
separation---information content, local geometry, and optimization
pathway---is the essential diagnostic structure of the workflow-limited
failure.

\subsection{Transfer to field applications}
\label{sec:disc-transfer}

The synthetic results support a transfer in logic rather than a transfer in
numerical thresholds. The exact noise and observation-density boundaries are
specific to the steady unit-square benchmark. Real aquifers may have pumping,
recharge, transient responses, geological boundaries, irregular observation
layouts, and independent measurements that change
identifiability~\cite{yeh2015flow, cardiff2013hydraulic}. The limitation also
follows the broader modeling principle that controlled models can be useful
diagnostic tools without being validated as direct field
surrogates~\cite{konikow1992validation}.

What is likely to transfer is the evaluation principle. Head fit should not be
equated with conductivity recovery. Reduced parameterizations should be matched
to the information content of the data. Direct deterministic baselines and
closed-loop validation should be reported before claiming PINN-based
conductivity inversion success. When sparse head-only data are the dominant
information source, unconstrained conductivity maps should be treated with
particular caution.

\subsection{Limitations and next steps}
\label{sec:disc-limit}

The current study is synthetic, steady, source-free, and uses regular boundary
conditions. The main noise matrix now includes a compact eight-truth FDM check, but
seed coverage is still concentrated in targeted audit cases rather than a full
factorial ensemble. The FIM diagnostic has been repeated across the compact
truth realizations for $\nobs=200$, at recovered points for two hard-failure
truths ($2$ and $3$, twelve paired cases), and across four observation
densities ($50$--$400$), but should still be extended to
truth/recovered-point combinations across densities.
The direct baseline is deterministic FDM-KLE least squares, not a complete
hydrogeologic inversion ecosystem~\cite{doherty2015pest, hill2007effective}.

The most important next extension is not to make the benchmark broader in
every possible direction, but to make the current diagnostic claim
statistically firmer. The 3-mode noise/density matrix has been expanded to eight
truth realizations for the direct FDM-KLE baseline (Section~4.3), and the
Jacobian diagnostic has been extended across observation layouts ($50$--$400$,
Section~4.5). A natural next stage would be to repeat the matched hybrid
comparison across the same expanded truth set and to extend the
recovered-point diagnostic across observation densities. More general hydraulic excitations, transient cases, and
Bayesian posterior analysis are valuable later extensions.

The current study deliberately isolates the staged A/B workflow because it
represents one practical way to separate head approximation from coefficient
recovery in PINN-KLE inverse studies~\cite{tartakovsky2020subsurface,
guo2023htpinn}. A joint end-to-end optimization that updates $h$ and $\bxi$
simultaneously may alleviate the stage-coupling failure mode (iv), but would
introduce its own convergence challenges in the combined landscape; evaluating
such alternatives is beyond the present diagnostic scope and would be a
valuable extension.

\subsection{Evidence-type labeling for identifiability claims}
\label{sec:disc-evidence}

The reduced-parameterization experiments demonstrate that identifiability
claims must be labeled by evidence type. Scalar effective conductivity is
structurally unidentifiable under the tested steady source-free setup because
changing a homogeneous scalar conductivity does not change the head solution.
Homogeneous axis-aligned anisotropy is also structurally unidentifiable under
the same left-right loading. A two-zone contrast becomes meaningful when an
absolute scale prior is imposed. Rotated anisotropy is numerically recoverable
in the tested homogeneous-tensor setting under multistart optimization, but local Hessian analysis reveals
shallow valleys.

These cases should not be collapsed into a single vague ``identifiability''
scale. They represent analytic structural arguments, numerical invariance
tests, optimization recoverability, local-curvature uncertainty, and
workflow-limited diagnostics. The PINN-KLE diagnostic framework in this paper
addresses the last two categories; it does not replace structural analysis or
geometric invariance arguments.

%% =====================================================================
\section{Conclusions}
\label{sec:conclusions}

This paper presents a controlled numerical diagnostic study of PINN-KLE
groundwater conductivity inversion under sparse steady head observations. The
main conclusions are:
\begin{enumerate}
\item In the tested staged PINN-KLE workflow, the head-observation loss
  provides no direct gradient pathway to KLE coefficients; updates are mediated
  through the PDE-residual branch in the implemented staged workflow.
\item No-truth round selection can identify the best available alternating
  round, but selecting the best round is not equivalent to conductivity
  recovery.
\item Direct FDM-KLE least-squares passes 24/24 cases across an expanded set of
  eight truth realizations ($\sigma_{\log K}=0.17$--$0.80$, three noise levels),
  showing that the sparse direct inverse problem still contains recoverable
  low-dimensional information under the tested 3-mode parameterization. In the
  matched three-truth seed audit, the PINN-hybrid workflow weak-passes only
  3/18 conductivity-recovery cases. This gap makes the moderate-noise hybrid
  failure workflow-limited in this setting.
\item Late-round hybrid degradation is Stage-B-led with frequent Stage A/B
  coupling.
\item Sparse-head Jacobian/FIM/Gauss-Newton diagnostics show linearized local
  observability for the direct map across observation densities from 50 to 400,
  with increasing conditioning burden as KLE mode count increases.
\item Recovery-point local conditioning remains comparable to true-point
  conditioning (condition-number ratio $0.919$--$1.010$, median $0.953$, across twelve paired
  cases), showing that the observed hybrid recovery failures are not caused
  by the workflow moving the coefficients into a locally unobservable region.
\end{enumerate}

The broader recommendation is methodological: PINN-based groundwater inverse
studies should distinguish head fitting, linearized local observability,
deterministic baseline recoverability, and closed-loop conductivity recovery.
Without this separation, a successful PINN head fit can be mistaken for a
successful hydraulic-conductivity inversion. This study focuses on diagnostic
information accessibility in a controlled benchmark, not on proposing a
field-ready inverse solver.

%% =====================================================================
\appendix
\section{Training and baseline configuration}
\label{sec:appendix-config}

This appendix summarizes the computational settings used in the reported
experiments. All values are extracted from the YAML configuration files and
source code in the working repository. The pilot experiments (Section~4.3,
pilot aggregate) use high-budget settings; the CG compact multi-truth matrix
and seed audit use medium-budget settings to keep the experimental matrix
computationally tractable. Both configurations are documented below.

\subsection*{A.1 PINN-hybrid workflow}

\begin{table}[!ht]
  \caption{Neural network and domain configuration.}
  \label{tab:config-nn}
  \centering\small
  \setlength{\tabcolsep}{3pt}
  \begin{tabular}{@{}p{0.38\columnwidth}p{0.54\columnwidth}@{}}
    \toprule
    Setting & Value \\
    \midrule
    Neural network depth & 5 hidden layers \\
    Neural network width & 128 neurons per layer \\
    Activation function & tanh \\
    Weight initialization & Xavier \\
    Hard Dirichlet BC & Yes (left/right enforced analytically) \\
    Fourier features & 8 frequencies, max 8.0, linear spacing \\
    KLE covariance kernel & Exponential, $\sigma^2\!=\!1.0$, $l_x\!=\!l_y\!=\!0.2$ \\
    FDM reference grid & $51 \times 51$ \\
    \bottomrule
  \end{tabular}
\end{table}

\subsection*{A.2 Stage-wise settings: high-budget (pilot) vs.\ medium-budget (CG compact)}

\begin{table}[!ht]
  \caption{Common optimizer and iteration settings.}
  \label{tab:config-budget}
  \centering\small
  \setlength{\tabcolsep}{3pt}
  \begin{tabular}{@{}p{0.31\columnwidth}p{0.28\columnwidth}p{0.30\columnwidth}@{}}
    \toprule
    Setting & High-budget (pilot) & Medium-budget (CG compact) \\
    \midrule
    Stage A iterations/round & 500 & 500 \\
    Stage B iterations/round & 1000 & 1000 \\
    PDE collocation points & 2000 & 1000 \\
    BC sampling points & 400 & 200 \\
    Alternating rounds (max) & 6--8 & 5 \\
    Closed-loop validation grid & $201 \times 201$ & $101 \times 101$ \\
    Observation seeds & 202, 303, 404 & 202, 303, 404 \\
    \bottomrule
  \end{tabular}
\end{table}

\begin{table}[!ht]
  \caption{Stage A (head fitting, $\bxi$ frozen). Common to both budgets.}
  \label{tab:config-stageA}
  \centering\small
  \setlength{\tabcolsep}{3pt}
  \begin{tabular}{@{}p{0.38\columnwidth}p{0.54\columnwidth}@{}}
    \toprule
    Setting & Value \\
    \midrule
    Optimizer & Adam \\
    Learning rate & $5 \times 10^{-4}$ \\
    Loss weights: PDE / BC / Data & 1.0 / 10--20 / 100.0 \\
    Weight decay / Gradient clip & $10^{-5}$ / 1.0 \\
    LR scheduler patience & Equal to iteration count (disabled) \\
    \bottomrule
  \end{tabular}
\end{table}

\begin{table}[!ht]
  \caption{Stage B (coefficient recovery, $h$ frozen). Common to both budgets.}
  \label{tab:config-stageB}
  \centering\small
  \setlength{\tabcolsep}{3pt}
  \begin{tabular}{@{}p{0.38\columnwidth}p{0.54\columnwidth}@{}}
    \toprule
    Setting & Value \\
    \midrule
    Optimizer & Adam \\
    Learning rate & $10^{-2}$ \\
    Coefficient initialization & From previous round (round 1: zero) \\
    $h$ smoothing $\sigma$ & 0.0 (no Gaussian smoothing) \\
    logK clamping & $\pm 3.0$ \\
    \bottomrule
  \end{tabular}
\end{table}

\subsection*{A.3 Direct FDM-KLE baseline}

The direct baseline optimizes the same low-dimensional KLE coefficients as the
PINN-hybrid workflow, but evaluates each candidate through the FDM forward
solver and minimizes the sparse-head residual directly. It is included as a
diagnostic comparator rather than as a claim of best-practice hydrogeologic
inversion.

\begin{table}[!ht]
  \caption{FDM-KLE baseline configuration.}
  \label{tab:config-baseline}
  \centering\small
  \setlength{\tabcolsep}{3pt}
  \begin{tabular}{@{}p{0.38\columnwidth}p{0.54\columnwidth}@{}}
    \toprule
    Setting & Value \\
    \midrule
    Optimizer & scipy.optimize.minimize (Powell) \\
    Objective & Sum of squared sparse-head residuals \\
    Forward model & FDM on $81 \times 81$ grid \\
    Max iterations & 80 \\
    Random starts & 2 \\
    Coefficient bounds & $\pm 3.0$ in $\bxi$ \\
    Closed-loop validation & FDM on $201 \times 201$ grid \\
    \bottomrule
  \end{tabular}
\end{table}

\subsection*{A.4 Jacobian/FIM diagnostic}

The local sensitivity diagnostic uses central finite differences of the
sparse-head observation map with respect to KLE coefficients. The reported FIM
and Gauss--Newton spectra are derived from the same Jacobian and are used only
as local linearized observability checks.

\begin{table}[!ht]
  \caption{Jacobian/FIM diagnostic configuration.}
  \label{tab:config-fim}
  \centering\small
  \setlength{\tabcolsep}{3pt}
  \begin{tabular}{@{}p{0.38\columnwidth}p{0.54\columnwidth}@{}}
    \toprule
    Setting & Value \\
    \midrule
    Finite-difference step & $\delta = 10^{-3}$ (absolute perturbation) \\
    Reference points & (a) true $\bxi$, (b) hybrid-recovered $\bxi$ \\
    Observation densities scanned & 50, 100, 200, 400 \\
    Noise level for FIM scan & $\sigma = 0.005$ \\
    \bottomrule
  \end{tabular}
\end{table}

%% =====================================================================
\backmatter

\section*{Acknowledgements}
The author thanks the China Geological Survey for institutional support.

%% ---- Declarations ----
\section*{Statements and Declarations}

\paragraph{Funding} This research did not receive any specific grant from funding agencies in the public, commercial, or not-for-profit sectors.

\paragraph{Competing interests} The author declares no competing interests.

\paragraph{Ethics approval} Not applicable. This study is a purely computational investigation using synthetic data and does not involve human participants or animals.

\paragraph{Consent to participate} Not applicable.

\paragraph{Consent for publication} Not applicable.

\paragraph{Data availability} The machine-readable summary tables and metadata
supporting the reported numerical results are supplied as Online Resource~2
for review, including the FDM-KLE baseline metrics, PINN-hybrid recovery
statistics, Jacobian/FIM diagnostics, and figure-level summary data. The raw
NumPy arrays for truth realizations, synthetic observation sets, and FDM
reference solutions are retained in the project output tree and are available
from the author upon reasonable request during review. A public
archive containing the review data package and the raw numerical arrays will be
deposited in Zenodo or HydroShare before publication, and the persistent
identifier will be added to the final article if the manuscript is accepted.

\paragraph{Code availability} The Python source code implementing the
finite-difference forward solver, PINN-KLE training workflow, FDM-KLE
least-squares baseline, Jacobian/FIM diagnostic scripts, and figure-generation
utilities is summarized in Online Resource~2 for review and is publicly
available at
\url{https://github.com/LancelYang/pinn-kle-groundwater-diagnostics}.
The training configurations (YAML) are self-documented in
Appendix~A, and the exact command-line invocations for each experiment are
recorded in {\tt runner\_command.json} files within each experiment output
directory. The computational environment uses Python~3.11+ with PyTorch and
NumPy/SciPy as the primary dependencies; dependency files are included in
Online Resource~2.

\paragraph{Author contribution} F.Y. conceived the study, designed the diagnostic framework, implemented the experiments, and wrote the manuscript.

\paragraph{Use of AI-assisted tools} During the preparation of this work, the
author used a large language model based coding assistant for code debugging,
workflow organization, and manuscript copy-editing. The author reviewed and
edited the generated content and takes full responsibility for the final
manuscript, code, and reported results.

%% ---- Bibliography ----
\bibliography{references}

\end{document}
